\section{Introduction}

\begin{frame}{}
    \tableofcontents[currentsection]
\end{frame}

\begin{frame}{Motivation}
    % The power of operads
    % lies in their ability to capture the essence of ``operations and their
    % symmetries'' in a purely combinatorial way, while simultaneously encoding
    % topological and algebraic information.

    \begin{thm}[Recognition principles for $E_n$-operads \cite{EO}]
        Let $A$ be an operad in the  category of compactly generated
        Hausdorff spaces. $A$ is an $\E_2$-operad if and only if $A(k)$ is connected and the collection of
        universal covering spaces $\{\Tilde{A}(k)\}_{k \in \N}$ forms a
        $B_{\infty}$-operad, i.e.:
        \begin{enumerate}[(1)]
            \item For all $k \in \N$, the space $\Tilde{A}(k)$ is contractible.
            \item The braid group  $B_k$ acts freely on $\Tilde{A}(k)$ for each
            $k \in \N$.
        \end{enumerate}
    \end{thm}

    \medskip
    \pause
    % \centering

    \textbf{Goal:} 
    
    Develop an analogous theory on braided operads as Lurie has given in 
    Higher Algebra (c.f. \cite[Section 2]{ha}) for symmetric operads.
\end{frame}

\begin{frame}{Main Results}
\begin{enumerate}
    \item Introduce {\em colored} braided operads.
    \item Show that they simultaneously generalize:
    \begin{itemize}
        \item unicolored braided operads.
        \item braided monoidal categories.
    \end{itemize}
    \item Introduce the $2$-category $\Brx$ {\em braided} pointed sets as a combinatorial variant of $\Ex$.
    \item Show that colored braided operads can be described as {\em operadic fibrations} over $\Brx$.
\end{enumerate}
\end{frame}